\documentclass[11pt]{article}
\usepackage[margin=1in]{geometry}
\usepackage{amsmath}
\usepackage{graphicx}
\usepackage{booktabs}
\usepackage{array}
\usepackage{longtable}
\usepackage[numbers,sort&compress]{natbib}
\usepackage[hidelinks]{hyperref}

\setlength{\emergencystretch}{2em}
\newcommand{\pending}[1]{\textbf{Results: pending} (produced by \nolinkurl{#1}).}
\newcolumntype{L}[1]{>{\raggedright\arraybackslash}p{#1}}
\newcommand{\sa}{$^\dagger$}
\newcommand{\ab}{$^\ddagger$}
% Every corpus count below is a macro generated by data/counts.py.
% Generated by data/counts.py from corpus.csv, reconciliation.csv and paper.tex. Do not edit.
\newcommand{\cBase}{10}
\newcommand{\cBaseAhead}{4}
\newcommand{\cBaseLLM}{2}
\newcommand{\cBaseLevel}{1}
\newcommand{\cBaseMixed}{3}
\newcommand{\cBasePeer}{5}
\newcommand{\cBaseRatePct}{97}
\newcommand{\cBaseReadOK}{8}
\newcommand{\cBib}{85}
\newcommand{\cConflicts}{11}
\newcommand{\cCtx}{21}
\newcommand{\cDirPred}{2}
\newcommand{\cEmp}{52}
\newcommand{\cEmpAbstract}{31}
\newcommand{\cEmpExt}{1}
\newcommand{\cEmpFull}{12}
\newcommand{\cEmpPartial}{9}
\newcommand{\cEmpPeer}{34}
\newcommand{\cEmpSingle}{3}
\newcommand{\cEmpUS}{22}
\newcommand{\cEmpYearTwentySix}{16}
\newcommand{\cIncluded}{56}
\newcommand{\cMultiUnitRef}{9}
\newcommand{\cNone}{31}
\newcommand{\cNoneReadOK}{7}
\newcommand{\cNotConflict}{4}
\newcommand{\cOOneN}{4}
\newcommand{\cOOnePeer}{1}
\newcommand{\cPFiveN}{3}
\newcommand{\cPFivePeer}{2}
\newcommand{\cPFourN}{4}
\newcommand{\cPFourPeer}{3}
\newcommand{\cPOneN}{4}
\newcommand{\cPOnePeer}{3}
\newcommand{\cPThreeN}{10}
\newcommand{\cPThreePeer}{5}
\newcommand{\cPairs}{15}
\newcommand{\cPartialPair}{1}
\newcommand{\cPlaceboYear}{10}
\newcommand{\cPos}{4}
\newcommand{\cReadOK}{21}
\newcommand{\cReconciled}{5}
\newcommand{\cRel}{8}
\newcommand{\cRows}{32}
\newcommand{\cRowsAbstract}{13}
\newcommand{\cRowsFull}{19}
\newcommand{\cSpread}{9}
\newcommand{\cSpreadDedup}{8}
\newcommand{\cSpreadDedupLower}{8}
\newcommand{\cSpreadLower}{9}
\newcommand{\cSpreadReadOK}{6}
\newcommand{\cUnexplained}{5}
\newcommand{\cUnseen}{4}
\newcommand{\cWithin}{3}
\newcommand{\EvidenceMapRows}{U1 marginal & 23 & 5 & 3 & 1 & 5 & 2 \\
U2 contrast / rank & 5 & 1 & 0 & 0 & 4 & 1 \\
U3 joint structure & 3 & 0 & 1 & 1 & 1 & 0 \\
U4 individual & 3 & 2 & 1 & 2 & 6 & 0 \\
U5 effect & 4 & 0 & 0 & 0 & 0 & 1 \\}


\title{Silicon Samples, 2022--2026: A Targeted Survey of How Language-Model Simulations of Human Populations Compare with Human Data}
\author{Real World Research\\\small Internal-review draft of 2026-10-11. Not peer reviewed.\\\small Named authors will replace the corporate byline before submission.}
\date{}

\begin{document}
\maketitle

\begin{abstract}
Since 2022, a growing literature has asked whether large language models (LLMs) conditioned on demographic or biographical descriptions can stand in for human survey respondents, experimental subjects, or whole populations. The answers look contradictory: some papers report ``algorithmic fidelity'' and strong correlations with real treatment effects, while others report that synthetic samples are no better than a uniform guess. We propose a taxonomy of four design choices that papers rarely separate: what the model is conditioned on, how a response is elicited, what unit the claim is about (population marginal, subgroup contrast, joint structure, individual, or effect), and what the result is compared against. Our targeted corpus has \cIncluded{} included works (\cEmp{} empirical studies and \cPos{} position papers), coded by one coder; \cEmpAbstract{} of the \cEmp{} empirical studies were coded from the abstract only, so every vote count below is reported with its full-text subset, and double coding is still required. As a consistency check, not a test, we examine \cConflicts{} cross-study pairs of results with opposite verdicts: an axis difference accounts for \cReconciled{} and part of another, \cUnexplained{} remain unexplained, and because \cBaseRatePct{}\% of all study pairs differ on some axis, the check has little power to fail. We tabulate reported results with the reference point their authors used, systematize six failure modes (variance collapse, demographic over-determination and stereotyping, Western default, acquiescence and asker sycophancy, contamination and recall, and instability across releases), and ledger the evidence for each proposed mitigation. Three patterns recur. Directions and rank orders agree with human data more often than levels, spreads or individual answers, though part of that gap reflects comparing scale-free correlations with absolute errors; joint structure is contested; and rankings can reverse with response format or release. In the studies that compared them, describing a distribution in one call matched human marginals better than aggregating persona calls, but in the smallest panels the gap is within finite-sample noise, and no study varies the elicitation method and the persona factorially. Of \cBase{} studies that fit a non-LLM baseline, the baseline is ahead in \cBaseAhead{}, level in \cBaseLevel{}, mixed in \cBaseMixed{}, and behind in \cBaseLLM{}, both of which give the LLM human data beyond a prompt. We close with open problems and a registration-ready protocol, not yet registered and with no results, giving one primary test for each proposition and for forward-dated validation.
\end{abstract}

\section{Introduction}

\citet{argyle2023} proposed that a language model conditioned on the backstories of real survey respondents produces ``silicon samples'' whose answer distributions track those of the corresponding human subgroups, a property they called algorithmic fidelity. The proposal arrived with obvious appeal. Surveys are slow and expensive, response rates are falling, and many populations of interest are hard to reach. If a model could stand in for a sample, social science, market research and policy analysis would change.

Four years later the evidence is large and conflicting. LLM predictions of the effects in 70 preregistered US survey experiments correlate strongly with the real effects \citep{ashokkumar2026}, and interview-grounded agents reach 83\% of 1,052 people's own test--retest consistency on held-out General Social Survey items \citep{park2024}. At the same time, ChatGPT's synthetic ANES respondents vary less than real ones, and 48\% of regression coefficients estimated from them differ significantly from the ANES estimates, with the sign flipping in 32\% of those \citep{bisbee2024}, although the same paper finds average scores that correspond closely to the ANES. Ten of 45 models score below a uniform guess on a 20-dataset benchmark \citep{hu2025simbench}; Pew Research Center's 2026 test finds its best configuration missing real panel answers by 11.2 to 14.7 points per wave \citep{pew2026}; and digital twins built from 500 answers per person correlate with those people's answers on new outcomes at an average $r = 0.20$ \citep{peng2025funhouse}.

Many of these results measure different estimands under different elicitation methods against different reference points. A paper reporting high correlation between simulated and real treatment effects across studies makes a claim about the direction and rank of effects; a paper reporting variance collapse makes a claim about the spread of a single item's answer distribution; a paper reporting low individual-level correlation makes a claim about people, not populations. This survey sorts the literature along those lines and checks how far that sorting goes. We use ``agree with human data'' rather than ``replicate'': most studies compare simulated answers with human answers, which is not replication of a finding.

\paragraph{Contributions.}
\begin{enumerate}
  \item A four-axis taxonomy of LLM population simulation (conditioning source, elicitation contract, unit of claim, reference point) with formal estimands, a coded matrix of the corpus on the four axes (Appendix~\ref{app:matrix}), and a consistency check of the taxonomy on \cConflicts{} cross-study conflicting pairs, with a base rate and a placebo axis that show how little power the check has (Sections~\ref{sec:taxonomy} and~\ref{sec:reconcile}).
  \item A table of reported results from 2022 to October 2026, in which every number carries the reference point its authors used or is marked as having none, with a pinpoint (version, table or section, subset) for each number in a committed extraction file (Section~\ref{sec:table}).
  \item A systematization of six failure modes, with measured size, proposed mechanism and strength of evidence (Section~\ref{sec:fail}), and a ledger of mitigations with the evidence for and against each (Section~\ref{sec:mitig}).
  \item Five propositions with falsification criteria and per-proposition evidence counts, and a registration-ready protocol that maps one primary test onto each proposition and onto forward-dated validation (Sections~\ref{sec:props} and~\ref{sec:protocol}). It is not yet registered, its results are pending, and no number in this paper comes from it.
\end{enumerate}

\paragraph{Disclosure.} The authors are building a commercial world-simulation product whose planned design is one this survey's conclusions favor; Section~\ref{sec:ethics} states the competing interest and how it shaped the paper.

\section{Related Work}
\label{sec:related}

\paragraph{Prior surveys.} \citet{mou2024}, published in this journal, survey LLM-agent social simulation and organize it by scale into individual, scenario and society simulation; their individual simulation, which mimics specific people or demographic groups, is our subject. \citet{vonderheyde2026review} systematically review 136 empirical studies of LLM use across the survey and public-opinion research process, before, during and after data collection, and report where LLMs succeed and fail; it supersedes the authors' earlier working paper. \citet{simmons2023} review LLMs as ``subpopulation representative models'', covering elicitation techniques and early implementations, and \citet{sarstedt2024} review silicon-sample comparisons for consumer and marketing research and find results that vary considerably by domain. We differ from these reviews in three ways. We restrict attention to studies that compare simulated survey or experimental responses with human data, which is one slice of Mou et al.'s individual simulation and one task in von der Heyde et al.'s pipeline. We code each study's unit of claim and reference point, which none of the reviews separates. And we ask, with stated limits, whether that coding accounts for specific conflicting results.

\paragraph{Frameworks with overlapping axes.} \citet{meister2024} vary question domain, steering method and distribution expression, which correspond to our item set, A1 and A2. CoMPosT characterizes simulations by context, model, persona and topic and measures caricature \citep{cheng2023compost}. Our A1 and A2 extend these with contracts that use no persona or a fine-tuned distribution, and A3 and A4 have no counterpart in either framework.

\paragraph{Statistical frameworks.} A separate literature treats model outputs as surrogates to be combined with a small human sample under valid inference: prediction-powered inference \citep{angelopoulos2023}, design-based supervised learning \citep{egami2023}, and its application to LLM-simulated subjects, the mixed subjects design \citep{broska2024}. \citet{ludwig2024} give an econometric framework in which prediction requires no training leakage and estimation requires a validation sample. Our conclusion in Section~\ref{sec:mitig} restates their thesis for population simulation.

\paragraph{Position papers.} Critical and qualified positions include \citet{dillion2023}, \citet{lin2025} on six fallacies of substitution, \citet{wu2025boundary} on variance-aware boundaries, \citet{agnew2024} on the ethics of substituting AI for participants, and \citet{anthis2025} for the method's promise under conditions. Broader discussions of AI in social science and its epistemic risks are \citet{grossmann2023}, \citet{bail2024} and \citet{messeri2024}.

\section{Scope and Method}
\label{sec:method}

\paragraph{Inclusion.} We include empirical studies, 2022 to 11 October 2026, in which an LLM produces answers to survey questions or experimental tasks intended to stand in for those of humans drawn from a defined population (survey respondents, experimental participants, annotators with stated demographics, or named individuals) and which compare those answers with human answers. We also include position papers that argue specifically about the validity of that substitution. We exclude work on LLMs as annotators of content without a target human population, on believability without comparison to real people, and on agent-based models scored only against aggregate behavioral series (information cascades, infections, employment). That last group is a distinct literature that already scores simulated dynamics against real data; we discuss it in Section~\ref{sec:open} (O2) and in the mitigation ledger rather than code it.

\paragraph{Corpus.} The bibliography has \cBib{} entries. Of these, \cIncluded{} are included works (\cEmp{} empirical studies, \cPos{} position papers); \cRel{} are related reviews and broad position pieces discussed in Section~\ref{sec:related}; and \cCtx{} are context citations outside the inclusion criteria, used for mechanisms or methods (for example contamination probes \citep{golchin2023}, forecasting benchmarks \citep{karger2024forecastbench}, archetype-based and social-media agent-based models \citep{chopra2024archetypes,yang2024oasis}, and believability without human comparison \citep{park2023generative}). The role of every entry, with the reason, is in \nolinkurl{data/corpus.csv}. Of the \cEmp{} empirical studies, \cEmpPeer{} appear in a peer-reviewed venue according to Crossref, the publisher, the ACL Anthology, OpenReview or the arXiv record, and \cEmpExt{} more as a peer-reviewed extended abstract; \cEmpSingle{} are single-author preprints; \cEmpUS{} use only US data; and \cEmpYearTwentySix{} first appeared in 2026, many of them within three months of this draft.

\paragraph{Search and its limits.} The corpus was assembled from a 100-note internal research archive, each note citing pages its author opened, supplemented by targeted searches, by forward and backward citation from anchor papers \citep{argyle2023,bisbee2024,santurkar2023,park2024}, and by works that internal reviewers of earlier drafts identified. We did not log the number of records identified or screened, so we cannot report a PRISMA flow; a logged search is part of the protocol (Section~\ref{sec:extract}). Every work cited here was opened by us on 2026-10-11 (arXiv record or HTML full text, publisher page, or Crossref record). Claims that something has not been done (``no study'', ``we did not find'') mean that we did not find it in this corpus.

\paragraph{Reading depth.} Of the \cEmp{} empirical studies, \cEmpFull{} were read in full, \cEmpPartial{} were coded from the abstract plus a search of the full text for the comparisons we code (the baseline comparison, data source and panel size), and \cEmpAbstract{} were coded from the abstract only. Abstracts rarely state reference points or baselines, so codes from abstracts are likely to understate them. Every vote count below gives the number of studies in it that were read at least in part in full text.

\paragraph{Extraction and coding.} One coder extracted every number in Tables~\ref{tab:dist}--\ref{tab:indiv} and~\ref{tab:mitig} into \nolinkurl{data/extraction.csv}, with the version or venue, the table, section or abstract it came from, and the subset it covers. In the tables, \ab{} marks a row whose numbers we checked against the abstract only; \cRowsFull{} of the \cRows{} rows in Tables~\ref{tab:dist}--\ref{tab:indiv} were checked against the full text. The same coder assigned A1--A4 codes to each included empirical study (\nolinkurl{data/corpus.csv}; Appendix~\ref{app:matrix}) and listed the result pairs in Table~\ref{tab:reconcile} (\nolinkurl{data/reconciliation.csv}). Every count in this paper is a macro generated from those files and from this paper's tables by \nolinkurl{data/counts.py}, so the text cannot drift from the files. The coding has not been checked by a second coder.

\paragraph{Reading the numbers.} We quote numbers as reported. Where we derive a number from reported ones, we say ``our arithmetic''. We never pool effect sizes across papers, because metrics, datasets and elicitation methods differ. Cross-study statements are vote counts with explicit denominators (``$k$ of $n$ studies that measure X''), drawn from the coded files, and they inherit the selection and reading limits above.

\section{A Taxonomy of Population Simulation}
\label{sec:taxonomy}

\paragraph{A1. Conditioning source.} What the model is told about the person or group: (a) nothing (the model's default voice); (b) a group label, such as a country or party; (c) a demographic profile, typically 3 to 15 variables, synthetic or sampled from marginals; (d) a demographic profile drawn from a real respondent's record, the original silicon-sample design \citep{argyle2023}; (e) that respondent's answers to other items \citep{ku2026,lee2026openweight,toubia2025twin}; (f) a long qualitative interview \citep{park2024}. Moving down the list costs more human data per simulated person, and (e) and (f) are no longer synthetic in the sense of requiring no new data collection. Group-level demographic sampling \citep{sun2024random} and synthetic backstories \citep{moon2024anthology} are variants of (c).

\paragraph{A2. Elicitation contract.} How a response is read out: (i) one sampled answer per call (persona aggregation); (ii) token log-probabilities over options; (iii) a verbalized probability vector for one persona or group \citep{meister2024,zhang2025vs}; (iv) a single ``describe the population'' call with no persona \citep{jang2026,doudkin2026}; (v) a model fine-tuned to emit distributions \citep{suh2025subpop,cao2025}. Contracts (i) and (ii) can disagree with each other: first-token probabilities and text answers mismatch in over 60\% of cases in some settings \citep{wang2024myanswer}, and forcing a multiple-choice format changes answers relative to unconstrained ones \citep{rottger2024}.

\paragraph{A3. Unit of claim.} What the result is said to establish: (U1) a population marginal; (U2) a subgroup contrast or ranking; (U3) joint structure, such as item correlations or regression coefficients; (U4) an individual's answers; (U5) the direction or size of a treatment effect. Fidelity at one unit does not imply fidelity at another. \citet{wang2026ess} finds moderate country-level recovery with negligible individual-level recovery in the same data, and \citet{williams2026beyond} find that the method that best matches marginals is not the one that best matches correlation structure.

\paragraph{A4. Reference point.} What a number is compared against, so that a reader can tell how good a given distance is: a uniform or random-response benchmark; the human sampling floor (the distance expected from sampling alone, including a sample-size equivalence); a human test--retest ceiling; a non-LLM predictor fit to human data; or another human benchmark, such as forecasters, annotators, or the gap between two human groups used as a yardstick. ``None'' means the authors compare the simulated quantity with the human quantity and give no such benchmark; the human value of the target statistic itself (a human entropy compared with a simulated entropy) is the comparison, not a reference point. Under this definition \cNone{} of the \cEmp{} empirical studies report no reference point; among the \cReadOK{} studies read at least in part in full text, \cNoneReadOK{} do. In a review of representative studies, \citet{wu2025boundary} found that fewer than half explicitly assess behavioral variance.

\subsection{Estimands}
\label{sec:estimands}

Let $P$ be the human population of a country $c$, observed through a survey sample with design weights $w_r$. For item $j$ with categories $k = 1,\dots,K_j$, the target marginal is $P_{cj}(k) = \sum_r w_r \mathbf{1}[y_{rj}=k] / \sum_r w_r$ over valid responses, and $\hat P_{cj}$ is its survey estimate. The simulator is $Q_{cj}(k \mid s, m, \kappa)$, the answer distribution produced by model $m$ under conditioning $s$ (A1) and contract $\kappa$ (A2). For a population-level contract, (iii) with a group label or (iv), $Q_{cj}$ is the stated distribution. For a per-persona contract it is the equal-weight mixture over the $n$ simulated respondents, $Q_{cj} = n^{-1}\sum_{r=1}^{n} q_{rj}$, where $q_{rj}$ is the one-hot answer under (i) or the stated vector under (iii), and the $n$ personas are drawn with probability proportional to weight, so the mixture is self-weighting. The units of claim are functionals of $P$ and $Q$:
\begin{itemize}
  \item U1: the marginal $P_{cj}$, scored by a distance $d(Q_{cj}, \hat P_{cj})$.
  \item U2: a contrast $\mu_{gj} - \mu_{g'j}$ between subgroup or country means, or the rank order of $\{\mu_{gj}\}_g$.
  \item U3: an item correlation $\operatorname{corr}(y_j, y_{j'})$ or a regression coefficient $\beta$ of $y_j$ on covariates.
  \item U4: the individual answer $y_{rj}$, scored by accuracy or correlation across respondents.
  \item U5: a treatment effect $\tau = \mathrm{E}[y \mid T=1] - \mathrm{E}[y \mid T=0]$.
\end{itemize}
For U1 we use the normalized score
\begin{equation}
S = \frac{d(U, \hat P) - d(Q, \hat P)}{d(U, \hat P) - d_{\text{floor}}}, \qquad d_{\text{floor}} = \mathrm{E}\left[d(\hat P^{*}, \hat P)\right],
\end{equation}
where $U$ is the uniform distribution and $\hat P^{*}$ is a design-based bootstrap replicate of the human estimate. Because $\hat P^{*}$ scatters around $\hat P$ as $\hat P$ scatters around $P$, $d_{\text{floor}}$ estimates the expected distance of the observed estimate from the true population. So $S = 0$ is uniform and $S = 1$ is as close to $\hat P$ as the true population is expected to be; a second independent sample of the same design sits further away (about $\sqrt 2$ times, for a distance that scales like a standard error) and scores below 1. $S$ can exceed 1, is unbounded below, and is unstable when $d(U, \hat P)$ is close to $d_{\text{floor}}$, that is, for items whose human distribution is near uniform, which include the contested items of most interest. The protocol therefore reports raw distances as primary for contrasts on the same items, aggregates $S$ as a ratio of sums, and excludes items with a small denominator from $S$ (Section~\ref{sec:protocol}). SimBench's score \citep{hu2025simbench} instead sets 100 at $d(Q,\hat P)=0$, an exact match to the observed finite-sample distribution, which no simulator can be expected to reach; its best score of 40.80 is a mean over items of the share of the distance from uniform that is closed, against an unreachable target. Applied to individual accuracy, the same normalization turns 71.72\% \citep{toubia2025twin}, against a 59.17\% random baseline and an 81.72\% test--retest ceiling, into $(71.72 - 59.17)/(81.72 - 59.17) \approx 0.56$ (our arithmetic).

\section{Reported Results}
\label{sec:table}

Tables~\ref{tab:dist}, \ref{tab:dist2} and~\ref{tab:indiv} collect the quantitative results; Tables~\ref{tab:dist} and~\ref{tab:dist2} cover U1--U3 and Table~\ref{tab:indiv} covers U4--U5. The A1/A2 column uses the codes of Section~\ref{sec:taxonomy}; ``?'' means we could not determine the code from what we read. \sa{} marks a single-author preprint; \ab{} marks numbers checked against the abstract only. Table~\ref{tab:map} counts included empirical studies by unit of claim and reference point.

\begin{table}[p]
\centering
\scriptsize
\caption{Population- and subgroup-level results (U1--U3). TV: total variation distance; W1: Wasserstein-1; MAE: mean absolute error in percentage points; SD: standard deviation. Pinpoints for every number are in \nolinkurl{data/extraction.csv}.}
\label{tab:dist}
\begin{tabular}{@{}L{1.9cm}L{2.7cm}L{5.4cm}L{0.9cm}L{3.3cm}@{}}
\toprule
Study & Data / models & Reported result & A1/A2 & A4 reference point \\
\midrule
\citet{bisbee2024} & ANES feeling thermometers; ChatGPT & Means ``correspond closely'' to ANES; less variation; 48\% of coefficients differ significantly, of which 32\% flip sign; results shift over 3 months. Power example (affective polarization): SD 16.1 vs.\ ANES 31.4, effect 12.5 vs.\ 7.8 & c / i & None \\
\citet{santurkar2023}\ab & Pew ATP (OpinionQA), 60 US groups & Misalignment on par with Democrat--Republican divide on climate; persists after steering & a,b / ii & Other human: Democrat--Republican divide \\
\citet{dominguez2024}\ab & US ACS; 43 models & With answer order randomized, responses trend to uniform regardless of size & a,c / ii & Uniform \\
\citet{meister2024} & OpinionQA, GlobalOpinionQA, NYT books & Table 1 (all three datasets): verbalized JSON TV 0.226 (best); 30-sample sequences 0.278--0.328; log-probs 0.455--0.550. Table 2 (OpinionQA + NYT only): GPT-4 verbalized 0.204 & b / i,ii,iii & Uniform 0.363 (Table 1); human annotators 0.250 (Table 2) \\
\citet{hu2025simbench} & 20 datasets, 130 countries; 45 LLMs (v4) & Best 40.80/100; ten models below 0; group conditioning lowers score (religiosity $-9.91$) & a,b / ii,iii & 0 = uniform, 100 = exact match \\
\citet{suh2025subpop} & OpinionQA; 4 base models & W1 zero-shot 0.138--0.209; fine-tuned 0.094--0.106 & b / v & Human floor 0.031; uniform 0.178 \\
\citet{zhao2024wvb}\ab & WVS, 94,728 people & Share of questions with W1 $<0.2$: 11.1\% (Alpaca-7B) to 75.0\% (GPT-3.5) & c / i & None \\
\citet{tao2024} & IVS, 107 countries; 5 OpenAI models & GPT-4o cultural distance 2.42 $\to$ 1.57 with country prompt; improves 71--81\% of countries for GPT-4, GPT-4-turbo and GPT-4o & a,b / ? & None \\
\citet{atari2023}\ab & Cross-cultural data and cognitive tasks & Similarity to humans falls with distance from WEIRD populations ($r=-.70$) & a / ? & None \\
\citet{ozkan2026}\sa & 2,414 WVS respondents (Turkey-first) & Entropy 1.46 $\to$ 0.77 nats; TVD 0.44; 85\% of 20 scenario$\times$seed units collapsed; Verbalized Sampling SD ratio 0.40--0.56 $\to$ 1.26--1.37 & d / i,iii & None \\
\bottomrule
\end{tabular}
\end{table}

\begin{table}[p]
\centering
\scriptsize
\caption{Population- and subgroup-level results (U1--U3), continued.}
\label{tab:dist2}
\begin{tabular}{@{}L{1.9cm}L{2.7cm}L{5.4cm}L{0.9cm}L{3.3cm}@{}}
\toprule
Study & Data / models & Reported result & A1/A2 & A4 reference point \\
\midrule
\citet{lee2026openweight} & 2,058-person US panel (cites Twin-2K-500); 3 open families & Null persona 1.24 bits, 60\% midpoint; 2.01 bits conditioned on Big Five answers (human 2.20); simulated cross-instrument correlations track human ones, $r=0.70$--$0.73$ over 139 pairs & a,e / i & Human $n{=}40$ subsample as sampling reference \\
\citet{doudkin2026}\sa & 6 batteries, 3 countries; panels up to 500 & Persona panels MAE 12.39; no-persona query 6.27, wins 9/9; probability vectors 5.34, beat query 4/9; 66/128 option slots empty & a,c / i,iii,iv & National table: 7.00 vs.\ 7.47--7.65 on 3 unpublished cells \\
\citet{jang2026} & 100 OpinionQA (Pew ATP) items; 100 personas per item & Describe-the-distribution call TV 0.22 vs.\ persona aggregation 0.46; same persona, same answer on $>$half of items & c / i,iv & None \\
\citet{kim2026korean}\ab & Korea Media Panel; about 8,000 personas per model & MAE 14--19 pp; binary item-mean $r=0.70$--$0.91$; calibrated on 30\% holdout 4.9/4.4 pp & c / i & Direct estimate from same holdout 3.6 pp \\
\citet{wang2026ess}\sa & ESS R11, 30 countries & Country name lifts median item correlation $-0.03 \to 0.52$; individual recovery negligible & b,c / i? & Neighbor-country average: error 0.056 vs.\ 0.180, better on 41/42 items; rank correlation 0.55 vs.\ 0.44 \\
\citet{krsteski2025} & Two panel surveys & Synthesis-only bias 24--86\%; $<$5\% with rectification & c / i,v & Human-only estimate \\
\citet{lukauskas2026} & Lithuania, $n=263$; 37 models & Acquiescence shift $+0.84$ SD; mean inter-LLM PSS 0.73; synthetic-trained regressors $R^2=-0.18$ on humans vs.\ 0.28; LLMs reproduce the direction of psychometric relationships & c / i & Copula 0.69 vs.\ best LLM 0.71 PSS; copula ahead on sample-driven components; human ceiling 0.825 \\
\citet{peng2025funhouse} & 19 studies, 164 outcomes & Twin SD below human in 154/164; means differ in 105/164; mean gap 0.352 SD & a,c,e / i & MSE of the mean about that of a 10-person human sample mean \\
\citet{tornberg2026}\ab & 6 frontier models; 30,990 responses & Conservative asker cue: share of items closer to Democrats falls 28--62 pp; rightward shift 8.0$\times$ leftward & a / i & None (no-cue contrast) \\
\citet{chen2026synthetic} & GSS + WVS; 4 models & Between-segment gaps inflated 2--4$\times$ & c / i & Non-LLM baselines \\
\citet{pew2026} & Pew ATP waves 185/190/192, fielded Jan--Apr 2026 & Opus 4.6 configuration: 11.3 (methodology page 11.4), 14.7, 11.2 pts avg.\ abs.\ error by wave. GPT-5.1 13.3 (wave 185 subsample) and 16 (standard profile, full sample), not comparable & d / i & None \\
\citet{kozlowski2024}\ab & US liberals and conservatives; LLM trained on text through 2019; COVID-19 items & Reproduce observed partisan differences in 84\% of cases & b / ? & Chance \\
\citet{boelaert2025}\ab & Opinion polls; generative LLMs & Cannot replace respondents; strong bias and low variance on each topic; bias varies randomly between topics & ? / ? & None \\
\bottomrule
\end{tabular}
\end{table}

\begin{table}[!htbp]
\centering
\scriptsize
\caption{Individual-level (U4) and effect-level (U5) results.}
\label{tab:indiv}
\begin{tabular}{@{}L{1.9cm}L{2.7cm}L{5.4cm}L{0.9cm}L{3.3cm}@{}}
\toprule
Study & Data / models & Reported result & A1/A2 & A4 reference point \\
\midrule
\citet{park2024}\ab & 1,052 US adults & Held-out GSS: interview 83\%, survey 82\%, both 86\% of test--retest; demographics-only 74\% & c,e,f / ? & Test--retest ceiling \\
\citet{toubia2025twin} & Twin-2K-500, 2,058 US adults & Main configuration (text persona, GPT-4.1-mini) 71.72\%; best in Table 2 71.92\% (JSON persona, GPT-4.1) & e / i & Random 59.17\%; test--retest 81.72\% \\
\citet{peng2025funhouse} & 19 studies, 164 outcomes & Twin--human $r=0.197$ on new outcomes; demographics-only 0.145; empty persona 0.080 & a,c,e / i & Random responses 0.001 \\
\citet{ku2026} & TEDS 2024, Taiwan; 27B--120B & 52\% on unseen items; 67\% (partisan) to 23\% (sovereignty) & e / i & Random forest on same population 58.3\% vs.\ best LLM 52.1\% \\
\citet{chen2026synthetic} & GSS + WVS; 4 models & GSS: best LLM ties a demographic lookup (0.589) and trails logistic regression (0.622); WVS: every model 11--22 points below the baseline & c / i & Non-LLM baselines \\
\citet{ashokkumar2026}\ab & 70 US survey experiments, 469 effects, 119,330 people; GPT-4 & Strong correlation with real effects; remained high for studies not posted by the training cutoff; effect sizes overestimated; 15 megastudies: lower & a / ? & Similar to pooled human forecasters; megastudies: comparable to pooled expert forecasters \\
\citet{aher2023}\ab & 4 classic experiments & 3 of 4 replicated; ``hyper-accuracy distortion'' on wisdom of crowds & c / i & None (original findings) \\
\citet{park2023diminished}\ab & 14 Many Labs 2 studies; GPT-3.5 & 37.5\% of 8 analyzable studies replicated; 6 unanalyzable from near-zero variance & a / i & None (original findings) \\
\citet{binz2025centaur}\ab & Psych-101: 160 experiments, $>$60,000 people & Captures held-out participants' behavior better than existing cognitive models; generalizes to new cover stories and domains & e / v & Cognitive models \\
\bottomrule
\end{tabular}
\end{table}

\begin{table}[!htbp]
\centering
\small
\caption{Evidence map: included empirical studies by unit of claim (A3) and reference point (A4), from \nolinkurl{data/corpus.csv}. A4 is coded per study, not per unit, so a study with several units is counted under each of its reference points for each unit; this can place a study in a cell it did not report. It affects the \cMultiUnitRef{} studies that report more than one unit and any reference point.}
\label{tab:map}
\begin{tabular}{@{}lrrrrrr@{}}
\toprule
Unit & None & Uniform & Floor & Ceiling & Non-LLM baseline & Other human \\
\midrule
\EvidenceMapRows
\bottomrule
\end{tabular}
\end{table}

\subsection{A consistency check on conflicting results}
\label{sec:reconcile}

Table~\ref{tab:reconcile} lists the \cPairs{} pairs of results that we first judged to conflict, and the axis, if any, that accounts for each. We then fixed a definition, after the pairs were chosen and so post hoc: two results conflict if they come from different studies whose authors draw opposite verdicts about whether the LLM approximates the human data. By that definition \cNotConflict{} of the \cPairs{} pairs are not conflicts: \cWithin{} are contrasts within one study that the authors themselves report as different units (pairs 6--8; Wang's title states the distinction), and in pair 10 both studies find the baseline ahead. Of the remaining \cConflicts{}, an axis difference accounts for \cReconciled{} and partly for a sixth (Argyle versus SimBench, where the populations also differ); \cUnexplained{} remain unexplained. Each unexplained pair turns on country, language, model release or item entropy, which the four axes do not separate. Country and language would need a fifth axis (target population); release and entropy point to instability that no static coding captures.

This is a consistency check, not a test, for three reasons. First, it can hardly fail: \cBaseRatePct{}\% of all pairs of included empirical studies differ on at least one axis, and a placebo axis, year of first public version, differs in \cPlaceboYear{} of the \cConflicts{} conflict pairs. Second, in only \cDirPred{} of the \cReconciled{} accounted-for pairs does the axis difference predict the direction of the disagreement from a proposition in Section~\ref{sec:props} (pairs 1 and 11, through P1 and P3); in the others the axis says why the results need not agree, not which should look better. Third, pairs 9 and 11 compare different tasks, and the pairs were chosen by the coder who built the taxonomy. Section~\ref{sec:extract} specifies the blind rebuild that would make it a test.

\begin{table}[!htbp]
\centering
\scriptsize
\caption{Result pairs first judged to conflict (\nolinkurl{data/reconciliation.csv}). Type: conflict = different studies, opposite verdicts; within = one study, distinction drawn by its authors; same = same verdict. Dir.: whether the axis difference predicts which result looks better, from P1--P5.}
\label{tab:reconcile}
\begin{tabular}{@{}rL{3.7cm}L{3.7cm}L{1.0cm}L{3.9cm}L{1.0cm}@{}}
\toprule
\# & Result A & Result B & Type & Accounting axis & Dir. \\
\midrule
1 & Effect predictions correlate strongly \citep{ashokkumar2026} & 32\% of differing coefficients flip sign \citep{bisbee2024} & conflict & A3: U5 across studies vs.\ U3 within one & yes (P1) \\
2 & 83--86\% of test--retest \citep{park2024} & $r=0.20$ on new outcomes \citep{peng2025funhouse} & conflict & A4: ceiling-normalized vs.\ raw; also item scope & no \\
3 & Algorithmic fidelity \citep{argyle2023} & Ten models below uniform \citep{hu2025simbench} & conflict & Partly A1 and A4; populations differ & no \\
4 & Verbalized distributions beat uniform \citep{meister2024} & Trend to uniform once order randomized \citep{dominguez2024} & conflict & A2: verbalized vector vs.\ log-probs & no \\
5 & Means close to ANES \citep{bisbee2024} & 11.2--14.7 pt errors \citep{pew2026} & conflict & Unexplained (both US, U1) & --- \\
6 & Item-mean $r$ 0.70--0.91 & MAE 14--19 pp (both \citealp{kim2026korean}) & within & A3: U2 ordering vs.\ U1 level & --- \\
7 & Country correlation 0.52 & Negligible individual recovery (both \citealp{wang2026ess}) & within & A3: U2 vs.\ U4 & --- \\
8 & Fine-tuning best on marginals & Persona prompting marginally better on correlations (both \citealp{williams2026beyond}) & within & A3: U1 vs.\ U3 & --- \\
9 & No-persona query beats panels \citep{doudkin2026} & Richer conditioning raises accuracy \citep{park2024} & conflict & A3: U1 vs.\ U4 (different tasks) & no \\
10 & Within 6 pp of a random forest \citep{ku2026} & No LLM beats strongest baseline \citep{chen2026synthetic} & same & --- & --- \\
11 & Similar to pooled forecasters \citep{ashokkumar2026} & Copula ahead on sample-driven components \citep{lukauskas2026} & conflict & A4: forecasters vs.\ fitted model; A3 (different tasks) & yes (P1, P3) \\
12 & Moderately faithful means \citep{lam2026fast} & Inaccurate vote choice in Germany \citep{vonderheyde2025} & conflict & Unexplained (country) & --- \\
13 & Fidelity scales log-linearly with size \citep{hu2025simbench} & Newest Llama has lowest sign agreement \citep{lee2026openweight} & conflict & Unexplained (release; 14-pair panel) & --- \\
14 & Translation does not move answers \citep{durmus2023} & Questionnaire language steers best \citep{liu2025global} & conflict & Unexplained (language) & --- \\
15 & Instruct worse than base in every family \citep{jang2026} & Instruct beats base above 10B \citep{hu2025simbench} & conflict & Unexplained by axes; item entropy moderates & --- \\
\bottomrule
\end{tabular}
\end{table}

\section{Where Simulations Agree with Human Data}
\label{sec:replicate}

\paragraph{Effect directions and rank orders of means agree best, but rankings are fragile and the comparison is partly metric-driven.} Across 70 survey experiments, predicted effects correlate strongly with real ones, at an accuracy similar to pooled human forecasters, and correlations remained high for studies not published by the model's training cutoff \citep{ashokkumar2026}. A model trained on text through 2019 reproduced observed partisan differences in COVID-19 attitudes in 84\% of cases \citep{kozlowski2024}, and item-mean ordering in a non-English panel comes out ($r=0.70$--$0.91$; \citealp{kim2026korean}). The qualitative pattern of classic experiments also comes out \citep{horton2023,aher2023}, but those experiments and their published results are almost certainly in pretraining data, so they are weak evidence of simulation; the post-cutoff tests of Ashokkumar et al.\ and Kozlowski et al.\ are stronger. Rankings depend on the contract and the release: country rankings in the ESS run backwards unless scale endpoints are named in words \citep{wang2026ess}, and the newest of three Llama releases had 71.4\% sign agreement against 92.9\% for the two earlier ones, on a 14-pair panel the authors call indicative only \citep{lee2026openweight}. Magnitudes are biased even where directions hold: effect sizes are overestimated \citep{ashokkumar2026} and between-segment gaps are inflated two- to fourfold \citep{chen2026synthetic}.

Part of the gap between ``ranks'' and ``levels'' comes from the metrics. Rank and direction fidelity are reported as correlations, which are scale-free and grow when items or countries differ a lot in their means; level fidelity is reported as absolute error in percentage points. A trivial predictor, such as the pooled or previous-wave marginal, can reach high item-mean correlations while missing levels. In the one study that compared rankings with such a predictor, an average of neighboring countries ranked countries about as well as the best model condition and set levels better \citep{wang2026ess}. The cross-study pattern is therefore a hypothesis; H4 tests it on a common scale against a carry-forward comparator.

\paragraph{US levels: contested.} Some US studies find levels close to human ones: ChatGPT's average thermometer scores correspond closely to the ANES \citep{bisbee2024}, a Singapore study finds moderately faithful means \citep{lam2026fast}, and group-level demographic sampling produces distributions ``remarkably similar'' to US polls \citep{sun2024random}. Others find large US level errors: 11.2--14.7 points per wave on Pew's panel \citep{pew2026}, and means that differ significantly in 105 of 164 outcomes \citep{peng2025funhouse}. The taxonomy does not explain this disagreement (pair 5).

\paragraph{Joint structure: contested.} Against: 48\% of ANES regression coefficients differ significantly, and 32\% of those flip sign \citep{bisbee2024}; the fine-tuning method that best matches marginals is not the one that best matches correlations \citep{williams2026beyond}; and a Gaussian copula fit to the human data beats every LLM on the correlation and reliability components of a psychometric score \citep{lukauskas2026}. For: when personas are conditioned on real respondents' answers to one instrument, simulated correlations with a second instrument track the human correlations at $r = 0.70$--$0.73$ across 139 pairs, mainly through correct sign rather than magnitude \citep{lee2026openweight}; LLMs reproduce the qualitative direction of human psychometric relationships \citep{lukauskas2026}; persona prompting is marginally better on correlations than fine-tuning \citep{williams2026beyond}; and algorithmic fidelity was proposed partly as agreement in patterns of relationships between variables \citep{argyle2023}. On this evidence the signs of associations often agree, their magnitudes often do not, and conditioning on real answers (A1 e) is where agreement has been shown.

\paragraph{Individual answers agree when the model is given that individual's answers.} Agents conditioned on a person's interview or survey responses predict held-out items from the same domain well relative to that person's own reliability \citep{park2024,toubia2025twin}. Zero-shot cross-item prediction comes within 6 points of a random forest trained on the same population \citep{ku2026}, and LLMs fine-tuned on survey data can retrodict missing GSS opinions, though they do less well on opinions never asked \citep{kim2023augmented}. The gain over demographics alone is real but modest (74\% to 83--86\% of test--retest in \citealp{park2024}). On new outcomes, twins built from 500 answers correlate with people at $r = 0.20$ \citep{peng2025funhouse}. Peng et al.\ do not break accuracy down by distance from what the person revealed, so whether accuracy falls with topic distance is our inference, not their finding. These designs need human data for every simulated person, so they do not scale to synthetic populations.

\paragraph{Marginals agree only with the right contract, and do not certify simulation.} Population marginals can be matched moderately well when the model is asked to describe the distribution directly \citep{meister2024,jang2026}, when a supervised calibration map is fit \citep{kambhatla2025}, or after fine-tuning on survey distributions \citep{suh2025subpop,cao2025}. A no-persona query matches published marginals better than persona panels do, in 9 of 9 comparisons \citep{doudkin2026}, so marginal agreement is evidence about the contract, not about whether individual respondents are being simulated. Comparisons between a stated distribution and a sampled panel carry a finite-sample handicap for the panel: for $K = 4$ categories near uniform, the expected TV between an $n$-draw empirical distribution and its own expectation is about 0.13 at $n = 30$ and 0.07 at $n = 100$ (our arithmetic, normal approximation). Meister et al.'s sequence contract emits 30 samples, so its gap to the verbalized contract (0.278 against 0.226 TV) is within that noise; Jang et al.'s 100-persona panels trail the describe call by 0.24 TV, well beyond it; and Doudkin's panels of up to 500 respondents come with a sampling-floor analysis. We did not find a finite-sample correction in Meister et al.\ or Jang et al.

\paragraph{Spreads and non-Western levels do not agree.} In all \cSpread{} included studies that compare simulated with human spread under persona sampling (\cSpreadReadOK{} of them read at least in part in full text), the simulated spread is lower (Section~\ref{sec:fail}); the one over-dispersed result comes from a different contract, Verbalized Sampling \citep{ozkan2026}. Level errors grow with distance from the populations best represented in training data \citep{atari2023,tao2024,liu2025global}, and vote-choice predictions in Germany were inaccurate and biased toward the Greens and the Left \citep{vonderheyde2025}.

\section{Six Failure Modes}
\label{sec:fail}

\subsection{Variance collapse}
Simulated respondents are more alike than real ones. Bisbee et al.\ report less variation than in the ANES \citep{bisbee2024}. In the power example in their Table 1 (affective polarization among partisans), the synthetic SD is 16.1 against 31.4 in the ANES; this is one quantity, and we do not generalize its ratio to other items. Twins had a smaller SD than their human counterparts in 154 of 164 outcomes \citep{peng2025funhouse}. Generative models answering opinion polls show low variance on each topic together with a strong bias that varies randomly from topic to topic, which \citet{boelaert2025} call machine bias. In WVS-grounded agents, 17 of 20 scenario$\times$seed units (85\%; our arithmetic from the reported rate) collapsed, meaning the agents' modal concentration exceeded the human modal mass by more than 0.15 \citep{ozkan2026}. Unconditioned personas put 60\% of answers on the scale midpoint \citep{lee2026openweight}, and early GPT-3.5 runs returned a single ``correct answer'' with near-zero variation on six of 14 Many Labs studies \citep{park2023diminished}. Silicon samples in Singapore understate opinion variance \citep{lam2026fast}. Persona populations also collapse into narrow behavioral modes in personality and moral-reasoning tasks \citep{xiao2026chameleon}, though that study has no human comparison.

\emph{Mechanism.} Two accounts have evidence. \citet{zhang2025vs} attribute post-training mode collapse to typicality bias in preference data. \citet{jang2026} show that instruction-tuned models are worse samplers than their own base models in every family they could compare, while the same models can describe the distribution accurately in one call. SimBench finds instead that instruction-tuned models beat their bases above 10B parameters, but that post-training helps on low-entropy (consensus) items and hurts on high-entropy (contested) ones \citep{hu2025simbench}. One reading that fits both, which is ours and not tested, is that opinion items are mostly contested and that Jang et al.\ measure contract (i), where collapse shows most directly.

\emph{Consequence.} Collapsed variance is not cosmetic. In Bisbee et al.'s power example, the synthetic inputs imply that 33 partisan respondents suffice, ``almost an order of magnitude'' fewer than the ANES inputs require \citep{bisbee2024}. That gap comes from a larger effect and a smaller SD together: required $n$ scales with $(\mathrm{SD}/\text{effect})^2$, which gives a ratio of about 9.8, of which the SD difference alone accounts for about 3.8 (our arithmetic). A regressor trained on synthetic respondents loses predictive validity on real ones \citep{lukauskas2026}. Collapse is not unique to LLMs: in Taiwanese data, supervised models also under-disperse \citep{ku2026}, which is a reason to compare against a fitted baseline rather than against an ideal.

\subsection{Demographic over-determination and stereotyping}
Models treat identity as more predictive than it is. Between-segment gaps are inflated two- to fourfold across GSS and WVS \citep{chen2026synthetic}; digital twins show insufficient individuation and stereotyping \citep{peng2025funhouse}; and conditioning on a group lowers SimBench scores relative to the unconditioned population baseline, most for religiosity \citep{hu2025simbench}. In subjective-annotation data, persona variables explain under 10\% of the variance in human labels, and persona prompting gives modest gains \citep{hu2024persona}; sociodemographic prompting does not consistently improve, and sometimes worsens, prediction of subgroup judgments \citep{sun2023sociodemo}, and its effect varies with model, size, dataset and prompt wording \citep{beck2024}. Simulations of political and marginalized groups are especially prone to caricature \citep{cheng2023compost}. \citet{wang2024flatten} argue analytically and show in human studies with 3,200 participants that LLMs both misportray and flatten identity groups. Counterfactual swaps show education effects (mean $|d|=0.56$) dwarfing role (0.18) and gender (0.12) \citep{lukauskas2026}.

\subsection{Western default}
Default answers resemble those of the USA and some European and South American countries \citep{durmus2023}, and resemble English-speaking and Protestant European countries across five OpenAI models \citep{tao2024}; similarity to humans declines with distance from WEIRD populations ($r=-.70$; \citealp{atari2023}). For GPT-4, GPT-4-turbo and GPT-4o, naming the country improved alignment for 71--81\% of 107 countries, and it left GPT-4o at a mean distance of 1.57 \citep{tao2024}. Models over-align with a few countries and under-align with most \citep{liu2025global}. Ten models stayed biased toward the Netherlands, Germany, the US and Japan across 11 prompt languages \citep{bulte2026}. Alignment tuning created disparities across countries in global opinion agreement \citep{ryan2024}, and multilingual ability did not reliably bring cultural alignment \citep{rystrom2025}. Evidence on prompt language conflicts: translating questions did not move answers toward speakers of that language in \citet{durmus2023}, while matching the questionnaire's language steered better than other methods in \citet{liu2025global}; \citet{bulte2026} found explicit cultural framing more effective than language. Outside the US, absolute errors are large even when rankings hold (14--19 points in Korea, \citealp{kim2026korean}). In the European Social Survey, adding the country name drove most of the country signal and richer profiles added no consistent gain \citep{wang2026ess}. Fine-tuning on country-level results from two global surveys outperforms prompting, though even the best models struggle on unseen questions \citep{cao2025}.

\subsection{Acquiescence and asker sycophancy}
Across 37 models, synthetic respondents shifted toward agreement by 0.84 SD \citep{lukauskas2026}; in the same study the models' response data were more similar to each other (mean inter-LLM PSS 0.73) than to the human data. Across nine models, LLMs generally fail to reproduce the direction and size of known human response biases in survey design, especially after RLHF, and react to perturbations that do not move humans \citep{tjuatja2024}. Answers also depend on who the model thinks is asking. With a conservative-Republican asker cue, the share of items on which six frontier models' answers were closer to Democrats fell by 28--62 percentage points, while a progressive cue produced little change; asked who the default asker is, models named an auditor, researcher or academic \citep{tornberg2026}. A simulated population therefore reflects the questioner as well as the persona.

\subsection{Contamination and recall}
Benchmarks built from WVS, Pew, GSS and ANES are public and old enough to be in pretraining data, as are the textbook experiments in U5 studies. General contamination probes exist \citep{golchin2023}, forecasting benchmarks avoid the problem by scoring only questions unresolved at submission \citep{karger2024forecastbench}, and \citet{ludwig2024} state the requirement of no training leakage for any prediction use. In population simulation the direct evidence is thin. \citet{ozkan2026} report a placebo-controlled memorization test and an election backtest suggesting that aggregate results survive while subgroup and individual claims are contaminated by recall; \citet{lukauskas2026} find no relation between a verbatim-recall probe and their leaderboard (rank correlation 0.00). \cUnseen{} included studies score models against human data the model is unlikely to have seen. \citet{ashokkumar2026} report that correlations remained high for studies not published or posted by the model's training cutoff, the closest published pre/post-cutoff test, for experimental effects. \citet{kozlowski2024} simulate 2020 COVID-19 attitudes with a model trained on text through 2019. Pew's waves were fielded January--April 2026, replicated within weeks, and published in September 2026 \citep{pew2026}; Pew does not report model cutoffs, so this is an out-of-training test of the answers unless a model's data extend past fieldwork. \citet{doudkin2026} use unpublished demographic cells. \citet{liu2025global} find models align best with the most recent survey wave, and \citet{kim2026korean} find an overestimate of generative-AI use consistent with temporal misalignment. We found no study that compares fidelity on survey waves released before and after a model's cutoff, with the same items and population, against a carry-forward baseline.

\subsection{Instability across releases and prompts}
The same prompt to ChatGPT gave significantly different results over a three-month period \citep{bisbee2024}. On a 14-pair panel, the newest of three Llama releases had the lowest sign agreement \citep{lee2026openweight}. Wording changes that alter meaning produce large human--model discrepancies, even for a model fine-tuned on psychological data \citep{schroder2025}. Option order and labeling biases dominate survey responses until randomized \citep{dominguez2024}, and forcing or relaxing a multiple-choice format changes answers \citep{rottger2024}. In the ESS, unlabeled scale endpoints led the model to rank countries backwards, and naming the endpoints in words stopped it \citep{wang2026ess}. A result for one model version and prompt format is a result about that version and format.

\section{Mitigations and Their Evidence}
\label{sec:mitig}

\begin{table}[!htbp]
\centering
\scriptsize
\caption{Mitigations, the failure they target, and the evidence. The first and ninth rows describe the design of the authors' planned product (Section~\ref{sec:ethics}).}
\label{tab:mitig}
\begin{tabular}{@{}L{2.8cm}L{1.9cm}L{5.2cm}L{4.4cm}@{}}
\toprule
Mitigation & Targets & Evidence for & Evidence against / caveat \\
\midrule
Describe the distribution in one call & Variance, cost & Verbalized JSON best TV \citep{meister2024}; TV 0.22 vs.\ 0.46 for 100-persona panels \citep{jang2026}; no-persona query beats panels 9/9 \citep{doudkin2026} & No per-person output; probability vectors beat the query 4/9 \citep{doudkin2026}; Meister's 30-sample panels are within finite-sample noise of the verbalized contract (our arithmetic); on unpublished cells no contract beats the national table \citep{doudkin2026} \\
Verbalized Sampling & Variance & Diversity 1.6--2.1$\times$ on creative writing \citep{zhang2025vs}; fidelity $+7$ to $+10$ points on a 0--100 scale \citep{ozkan2026} & Overshoots to SD ratio 1.26--1.37 \citep{ozkan2026} \\
Prompt perturbation per persona & Variance & Error $-21\%$ at no extra cost \citep{jang2026} & One paper \\
Condition on real answers or interviews & Variance, stereotyping, individual accuracy & Entropy 1.24 $\to$ 2.01 bits \citep{lee2026openweight}; 83--86\% of test--retest \citep{park2024}; backstories up to 18\% better distribution match \citep{moon2024anthology} & Needs human data per agent; $r=0.20$ on new outcomes \citep{peng2025funhouse} \\
Value-profile personas + calibration & Western default & Outperform generic, country and sociodemographic prompts on most model--country pairs, largest gains in under-represented countries \citep{abels2026} & 10 WVS countries; effect sizes given only in figures \\
Country / cultural prompting & Western default & Improves 71--81\% of countries for GPT-4, GPT-4-turbo and GPT-4o \citep{tao2024} & Can reflect stereotypes \citep{durmus2023}; no improvement for the other 19--29\% (our arithmetic) \\
Supervised calibration map & Levels & More consistent gains than uncalibrated and sociodemographic prompting across 15 LLMs, 3 datasets \citep{kambhatla2025} & Alignment falls for highly aligned groups (e.g.\ Aus/NZ); few-shot only in an appendix \\
Fine-tune on survey distributions & Levels & 32--46\% W1 reduction, unseen surveys and groups \citep{suh2025subpop}; outperforms prompting on two global surveys \citep{cao2025} & US training data in \citet{suh2025subpop}; W1 still about 3$\times$ the human floor (our arithmetic, 0.094/0.031); best models struggle on unseen questions \citep{cao2025} \\
Archetypes: query per group, sample agents & Cost & Unemployment MSE 24.59 vs.\ 41.05 heuristic and 56.98 LLM-as-agent; infections 95.17 vs.\ 2914.73 and 4311.70 (Table 1 of \citealp{chopra2024archetypes}) & Aggregate series for one city, not survey answers; positive correlation with labor data in 3 of 5 boroughs; variance not assessed \\
Rectification (prediction-powered) & Aggregate bias & Bias 24--86\% $\to$ $<$5\%; effective sample size up to 14\% higher \citep{krsteski2025}; valid inference by construction \citep{angelopoulos2023,egami2023,broska2024} & Validates aggregates, not agents; needs a human sample \\
Small human-calibrated panel & Levels & Ahead of a real-only shrinkage estimator below roughly 250--860 real responses, by at most 2.3 pp \citep{kim2026korean} & With a 30\% real holdout, the direct estimate is better (3.6 vs.\ 4.4--4.9 pp) \\
Larger / newer model & All & Fidelity scales log-linearly with size \citep{hu2025simbench} & Not with inference-time compute \citep{hu2025simbench}; no fix in \citet{chen2026synthetic,peng2025funhouse}; newer can be worse \citep{lee2026openweight} \\
Temperature & Variance & --- & MAE rises above default temperature for all 3 models (Table 2 of \citealp{lee2026openweight}); best twin at $T=0$ \citep{peng2025funhouse} \\
\bottomrule
\end{tabular}
\end{table}

Table~\ref{tab:mitig} lists the mitigations. Two conclusions follow, both with the caveat that their evidence is thin and that the first favors the authors' own design. First, the cheapest moves that reduced error in the studies that tried them change the elicitation contract (describe the distribution, perturb prompts) rather than the persona. Second, every method that reliably reduces level error uses real human data at the population level, through calibration, fine-tuning, or rectification \citep{kambhatla2025,suh2025subpop,krsteski2025}. In the one study that compared a calibrated synthetic panel with real data alone, the calibrated panel helped only when real responses numbered below roughly 250--860; with a 30\% real holdout, the real data alone were more accurate \citep{kim2026korean}. On current evidence, LLM simulation is a prior to be corrected by data, not a substitute for it, which is the position of the surrogate-inference literature \citep{angelopoulos2023,egami2023,broska2024,ludwig2024}. Position papers on both sides agree on the conditions: report variance with means, match claims to the unit validated, and validate per application \citep{wu2025boundary,anthis2025,schroder2025,lin2025}.

\section{Five Propositions}
\label{sec:props}

We condense the systematization into five propositions. Each states what would falsify it, and Table~\ref{tab:propmap} maps each to its current evidence and to a protocol test. Every test uses the three-way decision rule of Section~\ref{sec:protocol}.

\begin{description}
  \item[P1 (unit).] Fidelity holds at the unit validated: rank and direction fidelity (U2, U5) exceeds level fidelity (U1) on the same data, and exceeds spread, joint-structure and individual fidelity (U3, U4) \citep{wang2026ess,williams2026beyond,ashokkumar2026,kim2026korean}. Across studies this rests on correlations compared with absolute errors (Section~\ref{sec:replicate}), so it is a hypothesis. \emph{Refuted if}, on the same items and countries, the precision term of Lin's concordance does not exceed its accuracy term by the SESOI (H4, which tests the U1--U2 slice). Counter-evidence: US levels close to human in \citet{bisbee2024} and \citet{lam2026fast}; rankings reversed by format in \citet{wang2026ess}; joint-structure signs agree in \citet{lee2026openweight}.
  \item[P2 (contract).] Varying the elicitation contract changes fidelity more than varying the conditioning source. This is a conjecture: no study in the corpus varies both factorially. \citet{jang2026} and \citet{doudkin2026} compare contracts while removing the persona; \citet{park2024} varies conditioning only and finds it matters (74\% to 83--86\%); and probability vectors beat the no-persona query on 4 of 9 pairs \citep{doudkin2026}. \emph{Refuted if} the contract effect is not larger than the conditioning effect by the SESOI over the cells of H2, which exclude contract (iv).
  \item[P3 (baseline).] Against a non-LLM baseline given no more information, a prompted LLM is not ahead. Of the \cBase{} included studies that fit a non-LLM baseline (\cBaseReadOK{} read at least in part in full text), the baseline is ahead in \cBaseAhead{} \citep{chen2026synthetic,wang2026ess,ku2026,hu2024persona}, level in \cBaseLevel{} \citep{doudkin2026}, mixed in \cBaseMixed{} \citep{lukauskas2026,kim2026korean,kim2023augmented}, and behind in \cBaseLLM{} \citep{binz2025centaur,krsteski2025}. A code of ``ahead'' means the baseline is better on the authors' headline metrics and not worse on any; ``level'' means within the reported interval or described by the authors as a tie; ``mixed'' means better on some metrics or regimes and worse on others. Both LLM wins use human data beyond a prompt: a model fine-tuned on more than 60,000 people, and an LLM--human hybrid. Several baselines are given more information than the LLM: Ku's random forest and Hu and Collier's regression are fit on labeled data from the same population. The rule for each code is in \nolinkurl{data/corpus.csv}; \cBasePeer{} of the \cBase{} are peer reviewed. \emph{Refuted if} the LLM beats a baseline that uses no target-country labels on more country--item pairs than the SESOI allows (H3).
  \item[P4 (distance).] Error grows with cultural and linguistic distance from the populations best represented in training and preference data \citep{atari2023,tao2024,liu2025global,bulte2026}. Bult\'e and Rigouts Terryn measure Hofstede values and do not model a distance gradient, so their support is indirect. \emph{Refuted if} the partial correlation of per-country fidelity with distance is weaker than the SESOI (H5).
  \item[P5 (time).] Fidelity measured on one release does not transfer to the next, and fidelity measured before a model's cutoff overstates fidelity after it, so a validated simulator must be revalidated continuously \citep{bisbee2024,lee2026openweight,liu2025global}. Counter-evidence: correlations stayed high after the cutoff in \citet{ashokkumar2026}, and pre-pandemic text anticipated pandemic partisanship \citep{kozlowski2024}. \emph{Refuted if} both release-to-release change does not exceed within-version drift by the SESOI (H8) and post-cutoff degradation is below its SESOI (H6).
\end{description}

\begin{table}[!htbp]
\centering
\scriptsize
\caption{Propositions, current support, and the protocol test for each. ``Studies'' counts included empirical studies cited in support (keys in \nolinkurl{data/counts.py}); ``peer rev.'' counts those in a peer-reviewed venue. Each test is refuted when its interval excludes effects as large as its SESOI (Section~\ref{sec:protocol}).}
\label{tab:propmap}
\begin{tabular}{@{}L{1.5cm}L{1.3cm}L{1.0cm}L{4.4cm}L{2.0cm}L{3.2cm}@{}}
\toprule
Proposition & Studies & Peer rev. & Countries of supporting studies & Protocol test & Refuted if \\
\midrule
P1 unit & \cPOneN{} & \cPOnePeer{} & US, Korea, WVS countries, 30 ESS countries & H4 & Precision minus accuracy term $<0.10$ \\
P2 contract & 0 factorial & 0 & --- & H2 & Contract minus conditioning effect $<0.02$ TV \\
P3 baseline & \cPThreeN{} & \cPThreePeer{} & US, Taiwan, Lithuania, Korea, 30 ESS countries, 3 survey countries; annotation and cognitive tasks & H3 & LLM win share above 0.5 (H3 supported) \\
P4 distance & \cPFourN{} & \cPFourPeer{} & Up to 107 countries; 11 prompt languages & H5 & Partial $\rho$ weaker than $-0.4$ \\
P5 time & \cPFiveN{} & \cPFivePeer{} & US; WVS countries & H6, H8 & Both effects below SESOI \\
O1 forward & \cOOneN{} partial & \cOOnePeer{} & US experiments; US panel; US partisans; 3 survey countries & H9 & LLM advantage over carry-forward $<0.10$ \\
\bottomrule
\end{tabular}
\end{table}

\section{Open Problems}
\label{sec:open}

\paragraph{O1. Out-of-time validation.} Only \cUnseen{} included studies score models against human data the model is unlikely to have seen: post-cutoff experiments \citep{ashokkumar2026}, pandemic-era attitudes simulated by a pre-pandemic model \citep{kozlowski2024}, Pew's 2026 waves \citep{pew2026}, and unpublished cells \citep{doudkin2026}. Retrodiction of missing GSS years \citep{kim2023augmented} tests time in the other direction. Forward-dated evaluation, standard in forecasting \citep{karger2024forecastbench}, is rare for survey waves. We did not find a test in which simulated distributions for a survey wave were published before the wave's results were released and then scored against a carry-forward baseline.

\paragraph{O2. Dynamics.} The survey-response benchmarks we reviewed score static distributions. Outside our inclusion criteria, simulated dynamics have been scored against real data: an LLM-agent social-media simulator reproduced 198 real Twitter information cascades with a normalized RMSE of about 30\%, with cascade depth too low \citep{yang2024oasis}; LLM archetype agents forecast New York unemployment and infections with lower error than heuristic or per-agent LLM baselines \citep{chopra2024archetypes}; language models adapted to the media a subpopulation consumed predicted its survey answers on COVID-19 and consumer confidence \citep{chu2023}; and DARPA's SocialSim program scored pre-LLM simulators of online behavior against held-out data, with a baseline resampled from the period just before each window \citep{socialsim}. An LLM-agent opinion-dynamics study finds a strong bias toward scientifically accurate consensus unless confirmation bias is induced \citep{chuang2024}, but on our reading of the abstract it does not compare simulated trajectories with measured ones. What we did not find is a benchmark that scores the change in a population's survey answer distributions after real events against a persistence baseline, or against prediction markets.

\paragraph{O3. Global coverage of failure modes.} Variance collapse, acquiescence and sycophancy are measured mostly in US data, with single-country exceptions (Turkey, Lithuania, Korea, Taiwan, Singapore, Germany). Their size across many non-Western countries is unmeasured.

\paragraph{O4. Calibrated dispersion.} Methods either under-disperse (persona aggregation) or over-disperse (Verbalized Sampling). We did not find a dispersion calibration that is fit on held-out data and preserves joint structure.

\paragraph{O5. Joint structure.} Matching marginals is now routine; whether item correlations and regression coefficients agree is contested (Section~\ref{sec:replicate}): signs often agree when personas are conditioned on real answers \citep{lee2026openweight}, magnitudes and coefficients often do not \citep{bisbee2024,williams2026beyond,lukauskas2026}. Any downstream use that conditions or regresses on simulated data depends on it.

\paragraph{O6. Reporting standards.} We propose a minimal checklist: name the unit of claim; report variance and joint structure next to means; include uniform, marginal, sampling-floor and information-matched non-LLM baselines; report panel sizes and correct for finite-sample noise when comparing panels with stated distributions; pin model versions, decoding settings and dates; randomize option order and report refusal rates; report the asker framing; and split results by whether the human data were public before the model's cutoff.

\section{Registration-Ready Protocol}
\label{sec:protocol}

This section is a draft registration. It is not registered. Before any model call it will be frozen, committed to this repository with its analysis code and a synthetic-data dry run, and deposited on OSF with the repository commit hash; the OSF timestamp and commit hash will be the registration. Any later change is reported in a deviations section of the results paper, with the reason and date. \emph{No results exist yet; none are estimated here.} Running it is conditional on the data permissions in Section~\ref{sec:ethics}: if WVS has not given written permission within six months of registration, the protocol is not run, we report that, and any substitute source requires a new registration.

\paragraph{Data and eligibility.} World Values Survey wave 7 microdata is the primary source, with WVS wave 6 for previous-wave baselines, and Afrobarometer, Latinobar\'ometro and European Social Survey rounds added only where written permission for this use is obtained. GlobalOpinionQA is not used: its license (CC BY-NC-SA 4.0, per its dataset card) excludes commercial use and the authors are a company. An item is eligible if it has at least 2 substantive categories, at least 300 respondents with a valid answer in a country (an unweighted count), and is asked in at least 10 countries. ``Don't know'' and refusals are treated as missing in the human data and are not offered to the model; a sensitivity analysis keeps them as a category where the questionnaire offered them, and for H5 this is a pre-specified robustness check, because don't-know rates differ by region and so correlate with the predictor. Human distributions use the survey's documented design weights. From the eligible list, 40 target items are drawn at random, stratified by questionnaire module, with a committed seed. The item list, country list, and seeds are fixed in the registration.

\paragraph{Grid.} Contracts (A2): (i) one sampled answer per call; (i$'$) (i) with per-persona prompt perturbation \citep{jang2026}; (iii) a verbalized probability vector; (iv) one describe-the-population call. Conditioning (A1): (a) none, (b) country label, (d) a real respondent's demographic profile (age band, sex, education, urban/rural, country), (e) that respondent's answers to 20 items from modules other than the target item's module. Target items never appear in conditioning text. Per-persona cells, each with one call per simulated respondent: (i)$\times$\{b, d, e\}, (i$'$)$\times$d and (iii)$\times$\{d, e\}. Country-level cells: (iii)$\times$b and (iv)$\times$b; contract (iv) is defined only for population-level conditioning. (i)$\times$a runs once, without a country. The six cells \{(i), (iii)\}$\times$\{b, d, e\} form the factorial for H2. Option order is randomized per call. Three model families, one current version each (two consecutive versions where available, for H8), with pinned version identifiers and call dates; H6 adds older models (below).

\paragraph{Decoding and parsing.} Temperature 1.0 and top-$p$ 1.0 for every contract, with temperature 0 as a sensitivity analysis for (iii) and (iv) and the provider default for (i); where a provider fixes temperature for a model, the fixed value is recorded. Country-level cells are called 5 times and the stated vectors averaged; ``seed'' means the API seed parameter (values 1--5) where one exists and otherwise five independent calls. Reasoning budgets are set to the smallest configurable value for each family and recorded, with the provider default as a sensitivity analysis on (iv)$\times$b, because defaults differ across families and change over time. Maximum 512 output tokens. Answers are parsed strictly (option label or JSON); one retry with a format reminder is allowed; anything still unparseable or refused is excluded and counted. Refusal and parse-failure rates are reported by model, country and item, and any cell above 20\% is reported separately.

\paragraph{Sample sizes and calls.} Per-persona cells use $n = 300$ respondents per country, drawn with replacement with probability proportional to design weight and aggregated with equal weights, in 24 countries stratified by world region. Country-level cells run on every eligible WVS-7 country; H5 reuses the 24 countries already run. Table~\ref{tab:budget} gives the call counts implied by these choices (our arithmetic); across three families and one or two versions each, the total is roughly 6.2 to 12.4 million calls, plus the older H6 models. Cost per call is fixed from provider prices at registration.

\begin{table}[!htbp]
\centering
\scriptsize
\caption{Planned calls per model version (design arithmetic, not results). $C$ is the number of eligible WVS-7 countries, fixed at registration; the totals assume $C = 60$ and 30 country--item pairs for H6.}
\label{tab:budget}
\begin{tabular}{@{}L{2.4cm}L{6.6cm}r@{}}
\toprule
Block & Design & Calls \\
\midrule
H1--H4 & 40 items $\times$ 24 countries $\times$ (6 per-persona cells $\times$ 300 + 2 country cells $\times$ 5) + (i)$\times$a: 40 $\times$ 300 & 1,749,600 \\
H5 & 40 items $\times$ ($C - 24$) further countries $\times$ 2 country cells $\times$ 5 & 14,400 \\
H6 & $\geq$30 pairs $\times$ 2 waves $\times$ ((i)$\times$d 300 + (iv)$\times$b 5) + 240 cutoff-probe + 60 recall-probe calls & 18,600 \\
H7 & 40 items $\times$ 6 countries $\times$ 3 runs (2 framings + neutral replicate) $\times$ 300 & 216,000 \\
H8 & 200 canary items $\times$ 24 countries $\times$ (iv)$\times$b 5 calls $\times$ 3 dates & 72,000 \\
H9 & Target-wave items $\times$ countries $\times$ 305, fixed when the wave is chosen & --- \\
\midrule
Total (excl.\ H9) & & 2,070,600 \\
\bottomrule
\end{tabular}
\end{table}

\paragraph{Metrics and the score.} For nominal items, TV distance (primary) and Jensen--Shannon distance; for ordinal items, TV and W1 on category ranks scaled by $1/(K-1)$. The SD ratio and entropy are computed only for ordinal items with $K \geq 4$, coded $1,\dots,K$, which assumes interval spacing. Raw distances are primary for every contrast on the same items (H2, H3, H7, H8, H9). $S$ is used where floors differ between the things compared (H4, H5, H6), aggregated over items as a ratio of sums, $\sum_j [d_j(U) - d_j(Q)] / \sum_j [d_j(U) - d_{\text{floor},j}]$, and computed only over items with $d(U, \hat P) - d_{\text{floor}} \geq 0.05$ TV; excluded items are reported raw. $S$ is also reported with the pooled cross-country marginal in place of $U$, a stronger zero point.

\paragraph{Human sampling floor.} At full sample size, $d_{\text{floor}}$ is the mean distance between the weighted estimate and design-based bootstrap replicates, resampling primary sampling units within strata where the microdata identify them, and weighted respondents otherwise; in the second case the floor is inflated by the published design effect where one exists and flagged. Per-persona cells simulate 300 respondents, and even a perfect individual simulator would reproduce the 300 sampled respondents' own answers, not $\hat P$. Their floor is therefore the oracle panel: the mean distance between $\hat P$ and the empirical distribution of the sampled respondents' actual answers, over the same draws.

\paragraph{Finite-sample comparison.} A panel of $n$ draws has an empirical distribution whose TV from its own expectation is of order $\sqrt{K/n}$, while a stated distribution has no such noise. For every contrast between contracts (H2, H3) the primary comparison is at matched size: $n = 300$ answers are drawn from each stated distribution and scored like a panel. As a sensitivity analysis, each panel's distance is debiased by subtracting the mean distance between the panel and parametric-bootstrap resamples of itself; because TV is not additive, this over-corrects when the true distance is large, which favors panels.

\paragraph{Baselines.} Primary baselines use no labels from the target country and wave: (a) the pooled marginal across the other countries; (b) the neighbor-country average \citep{wang2026ess}, and the same country's WVS-6 marginal where the item was asked; (d) a multinomial logit on the same demographic variables, trained on respondents from the other countries; (e) a gradient-boosted classifier on the same 20 answered items, trained the same way. Logit and gradient-boosted baselines output probabilities, so they are compared at matched size by drawing one answer per sampled respondent. Secondary baselines are supervised on the target country's own respondents, disjoint from the evaluation respondents, as in the random-forest baseline of \citet{ku2026}; they are given strictly more information than the LLM, and we report a learning curve over training sizes 25, 50, 100, 250, 500 and 1,000 with the crossover size at which the baseline overtakes the LLM.

\paragraph{Decision rule.} Each primary has a parameter $\theta$, oriented so that the stated hypothesis predicts $\theta > 0$, a smallest effect size of interest $\delta$ (SESOI), and a 98.75\% interval, each of whose bounds is a one-sided test at $0.05/8$ (Bonferroni over the eight tested primaries). The hypothesis is \emph{refuted} if the upper bound is below $\delta$ (effects of interest are excluded, whatever the sign); \emph{supported} if the lower bound is above 0 and the upper bound is at least $\delta$; and \emph{inconclusive} otherwise. The three outcomes are exclusive and exhaustive. Primaries are pooled across the three model families, with family as a fixed effect and families weighted equally; per-family estimates are secondary. The synthetic dry run reports simulated power for each primary at its SESOI under the crossed design; any primary with power below 0.8 at the registered numbers of countries and items is either enlarged or relabeled as estimation-only before registration.

\paragraph{Hypotheses and primary tests.} Each hypothesis has one primary test; everything else is secondary and labeled exploratory.
\begin{description}
  \item[H1 (variance; estimation only).] The SD ratio, simulated to human, per model under (i)$\times$d, with an interval, next to the SD ratio of draws from the logit baseline (d), which separates model collapse from the information limit of demographics. Secondary: the same under (iii)$\times$d and (iv)$\times$b. No decision rule.
  \item[H2 (contract vs.\ conditioning; P2).] Outcome: matched-size TV in the six cells \{(i), (iii)\}$\times$\{b, d, e\}. With $m(\cdot)$ the marginal mean over the other factor, the contract effect is $|m(\text{iii}) - m(\text{i})|$ and the conditioning effect is $|m(e) - m(b)|$, the least against the most informative conditioning in the grid, fixed in advance so both effects are differences of two means. $\theta$ = contract effect minus conditioning effect; $\delta = 0.02$ TV. This tests P2 only over these cells; contract (iv), the strongest in the literature, is excluded because it is defined only with (b), and a secondary compares (iv)$\times$b with (iii)$\times$b.
  \item[H3 (baseline; P3).] $p$ is the share of country--item pairs where the (i)$\times$d cell has lower matched-size TV than the primary logit baseline (d); $\theta = p - 0.5$, the LLM's advantage, and $\delta = 0.10$. H3 is stated as the claim P3 denies: P3 is refuted if H3 is supported and supported if H3 is refuted. Secondary: every other cell against its baselines, the supervised learning curve, and U4 accuracy for (e).
  \item[H4 (unit; P1).] Under (iv)$\times$b over all eligible countries, for each item, Lin's concordance correlation between simulated and human country means factors into a precision term $\rho_j$ (Pearson correlation) and an accuracy term $C_{b,j}$ (shift in location and scale) \citep{lin1989}; $\theta$ = mean over items of $\rho_j - C_{b,j}$, $\delta = 0.10$. Both terms are on the same scale and equal 1 for a perfect simulator. Comparators: the same decomposition for the WVS-6 carry-forward predictor where available, and a country-permutation null for $\rho_j$. H4 tests the U1--U2 slice of P1. Secondary (U3, and O5): in per-persona cells, the correlation between the upper triangles of the simulated and human $40 \times 40$ item correlation matrices among the 300 sampled respondents, per country.
  \item[H5 (distance; P4).] Per-country $S$ under (iv)$\times$b against cultural $F_{ST}$ distance from the United States \citep{muthukrishna2020}, computed on WVS items disjoint from the 40 targets. $\theta = -\rho$, where $\rho$ is the partial Spearman correlation controlling for log human $n$, number of eligible items, questionnaire language (English or not) and survey mode; we do not control for mediators such as internet penetration. $\delta = 0.4$: Atari et al.\ report $r=-.70$, and a gradient below 0.4 would explain under a sixth of between-country variance in fidelity. At one-sided $\alpha = 0.05/8$, power to detect $\rho = -0.4$ is about 0.69 with 60 countries and 0.45 with 40, and 73 countries are needed for 0.80 (our arithmetic, Fisher $z$ with variance $1.06/(n-7)$); below 73 eligible countries H5 is reported as underpowered for its SESOI. Distance from the US is a proxy for representation in training data, and Atari et al.\ and Tao et al.\ frame the effect as distance from WEIRD or English-speaking Protestant countries; countries are not independent (Galton's problem), so a sensitivity analysis adds a world-region random effect.
  \item[H6 (post-cutoff degradation; P5).] For each country--item pair observed in an earlier and a later wave, degradation is $D = S_{\text{earlier}} - S_{\text{later}}$. Treated models have a training cutoff between the first public release of any aggregate from the two waves (toplines precede microdata); control models have a cutoff after both. $\theta$ = mean $D$ of treated minus mean $D$ of control models, so that positive means treated models lose more on the wave they have not seen; $\delta = 0.10$ in $S$. This identifies post-cutoff degradation, which combines staleness (no knowledge of real change between waves) with recall (memorized earlier toplines). Two secondaries separate them. A recall probe asks each model to complete published toplines of the earlier wave; recall predicts degradation concentrated in pairs with high recall accuracy, staleness predicts it regardless of recall. And items whose human distribution barely moved between waves (TV below the median) should show no degradation under either account, which checks for capability or format differences between waves. Cutoffs are taken from documentation and checked with a probe of 10 dated public events per month over 24 months, with the cutoff estimated as the change point of a logistic fit of accuracy on month. Waves: WVS-6 and WVS-7 unless permission for more frequent series (ESS rounds) is granted, so treated models need cutoffs between those releases, which in practice means older open-weight checkpoints; models whose pretraining corpora are public are preferred, because membership of the toplines can then be checked directly. The registration names the models. Feasibility gate: at least 3 models in each group and 30 country--item pairs; otherwise H6 is reported as infeasible, which we expect under a WVS-only plan. Within-model differencing removes level differences in capability but not their interaction with wave; the stable-item check addresses that.
  \item[H7 (asker framing; directional sycophancy).] Cell (i)$\times$d in the 6 per-persona countries with the highest valid response on the left--right self-placement item. The framing describes the asker, not the respondent: ``The person asking you this question places themselves at 2 (or 8) on a 0--10 left--right scale''; the respondent persona is unchanged. The reference distributions are those of human respondents in the same country who place themselves at 0--3 (left) or 7--10 (right). For each framing $f$, the directional shift is $\mathrm{TV}(Q_{\text{neutral}}, H_f) - \mathrm{TV}(Q_f, H_f)$, positive when framed answers move toward the asker's side. $\theta$ = the mean directional shift over both framings minus the same quantity computed with a second neutral run on a different date in place of $Q_f$; $\delta = 0.02$ TV, two points of probability mass, below a national poll's usual margin of error. Secondary: undirected TV between framed and neutral runs, and the asymmetry, right-framed minus left-framed shift.
  \item[H8 (release stability; P5).] A canary set of 200 items is sampled at random from the eligible list with a committed seed; the list is withheld, with a salted SHA-256 hash committed, until the earlier of results publication and 12 months after registration, when list and salt are published. Withholding protects only the selection against tuning on it during the window: the items and their human answers are public WVS data, every run sends them to the provider, and nothing stops WVS from entering training data. Each version is run on 3 dates at least 7 days apart; when a new release ships, old and new versions are run on the same dates. Release change is the mean over items of $|\mathrm{TV}_{\text{new}}(t) - \mathrm{TV}_{\text{old}}(t)|$ over shared dates $t$; drift is the mean over items of $|\mathrm{TV}_v(t) - \mathrm{TV}_v(t')|$ over date pairs and versions. Both are contrasts between two single-date runs. $\theta$ = release change minus drift, per pair of consecutive releases; $\delta = 0.02$ TV. If no new release in a family ships within six months of registration, H8 is reported as not run for that family.
  \item[H9 (forward-dated; O1).] For a survey wave whose aggregates are not yet public, chosen and named at registration, we publish in plaintext, before the first release, the simulated distributions under (iv)$\times$b and (i)$\times$d for every model and seed in a manifest fixed at registration, together with a carry-forward baseline from the previous wave and, where three or more prior waves exist, a linear-trend baseline. Every manifest entry is scored. A forecast needs no secrecy, so plaintext publication replaces hash commitments, which would allow committing many variants and revealing the best. $p$ is the share of country--item pairs where the (iv)$\times$b distribution has lower TV than carry-forward; $\theta = p - 0.5$, $\delta = 0.10$. Inference is conditional on the one wave.
\end{description}

\paragraph{Analysis.} All primary intervals come from a two-way crossed bootstrap over countries and items (the pigeonhole bootstrap, \citealp{owen2007}), because items and countries are crossed; respondents are not the resampling unit for model outputs, so within-cell persona noise enters only through the country--item residual. Every registered item is reported, including those where models beat or lose to baselines.

\paragraph{Systematic extraction.}
\label{sec:extract}
Before submission, the review itself is redone as a logged search. Sources: arXiv (cs.CL, cs.CY, cs.AI, econ.GN), the ACL Anthology, the Semantic Scholar API, SocArXiv, PsyArXiv and OSF Preprints, and the first 200 Google Scholar results per string, from 1 January 2022 to the registration date. Strings combine (``silicon sampl*'' OR ``synthetic respondent*'' OR ``simulat* survey respon*'' OR ``digital twin*'' OR ``algorithmic fidelity'' OR ``LLM persona*'') with (survey OR poll OR opinion OR experiment). Two coders screen titles and abstracts, then full texts, against Section~\ref{sec:method}, and record counts identified, screened and included. Every included study is coded from its full text. Both coders code every included study with the codebook in \nolinkurl{data/corpus.csv} (A1--A4 per unit, metric, value, reference value, model, country, data release date, model cutoff, panel size, finite-sample correction, peer-review status, and baseline outcome with the written rule) and extract numbers into \nolinkurl{data/extraction.csv}. Agreement on nominal codes must reach Krippendorff's $\alpha \geq 0.8$ \citep{hayes2007}; below that, the codebook is revised and the study recoded. A third coder adjudicates disagreements. For the conflict check, the definition of a conflict (different studies, same quantity class, opposite verdicts) and the directional predictions of P1--P5 are fixed first; each coder then lists conflicting pairs without seeing the axis codes; and the check reports the share of pairs an axis accounts for, the share where it predicts the direction, the same shares for a placebo axis (publication year), and the base rate over random pairs.

\paragraph{Results.} No result below exists yet.
\begin{itemize}
  \item H1--H5: \pending{bench/silicon-sok/results/h1-h5/}
  \item H6: \pending{bench/silicon-sok/results/h6-degradation/}
  \item H7: \pending{bench/silicon-sok/results/h7-asker-framing/}
  \item H8: \pending{bench/silicon-sok/results/h8-stability/}
  \item H9: \pending{bench/silicon-sok/results/h9-forward/}
  \item Double-coded extraction (required before submission): \pending{papers/silicon-samples-sok/data/extraction-double.csv}
\end{itemize}

\section{Limitations}
\label{sec:limits}

This is a targeted, not exhaustive, survey, and not a systematic review. We did not log the search, so we cannot report how many records were identified or screened, and a registered search could find studies that change the balance, especially in non-English venues. One coder extracted the numbers, coded the axes, and chose the result pairs; every vote count and the consistency check are single-coder judgments. Only \cEmpFull{} of the \cEmp{} empirical studies were read in full and \cEmpPartial{} more in part; \cEmpAbstract{} were coded from abstracts, which rarely report baselines, so the ``none'' count is likely too high. Of the \cRows{} rows in Tables~\ref{tab:dist}--\ref{tab:indiv}, \cRowsFull{} were checked against full texts and \cRowsAbstract{} against abstracts only (marked \ab). The tables mix metrics that are not comparable, which is why we count studies rather than pool effects. Several numbers are compared across studies with different panel sizes, and we have panel sizes for only three of them.

\emph{Dependence.} Studies are not independent replications. \citet{lee2026openweight} describe a 2,058-person national panel and cite Twin-2K-500 \citep{toubia2025twin}, a US sample of the same size, so we treat Lee et al.\ as US evidence. \citet{peng2025funhouse} and \citet{toubia2025twin} share authors and digital-twin infrastructure. Our vote counts do not adjust for this; treating these three as one source would change the spread count from \cSpreadLower{} of \cSpread{} to \cSpreadDedupLower{} of \cSpreadDedup{} and leave the other counts unchanged.

\emph{Publication and selection bias.} The corpus is skewed toward very recent work: \cEmpYearTwentySix{} of the \cEmp{} empirical studies first appeared in 2026, many as preprints posted within three months of this draft, and fast-moving preprints may be revised or withdrawn. Of the \cEmp{} empirical studies, \cEmpPeer{} appear in peer-reviewed venues; negative results appear in both peer-reviewed venues and preprints. Our selection, through an archive built for a product that depends on the answer, may favor either direction.

\emph{Scope.} Many results cover one country or one model family, and \cEmpUS{} of the \cEmp{} empirical studies use only US data. The frontier models in our sources are already being replaced, and P5 implies that specific numbers here will age quickly even where the patterns hold.

\section{Ethics Statement}
\label{sec:ethics}

\paragraph{Competing interests and authorship.} The authors are employees of Real World, a company building a commercial product that simulates populations, and this work was funded by that company. The planned product, Real World, is a daily simulation of about 100,000 synthetic agents in many countries, including the non-Western populations where this survey finds errors largest. Its population is designed to come from the GLOPOP-S synthetic population and published survey aggregates, not from census microdata. Each morning it is to correct segment beliefs against polls, with news and search trends as bounded forcing and prediction-market prices never assimilated; it commits hashed forecasts that are scored in public against the market price at the call time, and lets anyone ``fork'' the simulated world with an event and run a seven-day ensemble. The system and its evaluation are the subject of companion papers in preparation (\emph{Don't Ask a Model, Ask a World}, \emph{Receipts}, \emph{The Real World Benchmark} and others); none of their forecast results exists yet, and nothing in this survey depends on them. A first engine exists (version 0.1.0); it samples agents from illustrative marginals written for its repository rather than from GLOPOP-S, and asks the model once per segment for that segment's share answering yes, a form of the archetype design described next. Our internal design notes recommend querying the model once per archetype for a response distribution and sampling agents from it, the method of \citet{chopra2024archetypes} combined with contract (iii) or (iv). That design is what P2 and the first conclusion of Section~\ref{sec:mitig} favor, so we have a direct interest in both; we label P2 a conjecture and leave it to H2. Open problems O1 and O2 match the product's roadmap, and we chose to emphasize them knowing that. We have tried to counter the interest by leading with failure evidence and by stating P3, which cuts against the product. Named authors and a formal funding statement will replace the corporate byline before submission.

\paragraph{Data licensing.} We have not verified that any survey source permits this use. GlobalOpinionQA is CC BY-NC-SA 4.0 and is excluded. WVS downloads require agreeing to a non-redistribution data-use license whose commercial terms we have not read; Afrobarometer's usage policy states no license or commercial terms; we have not read the terms of Latinobar\'ometro or the ESS. A company validating its product is arguably commercial use. The protocol therefore runs on a source only after the archive gives written permission for this use, and a source that refuses or does not answer is dropped and reported.

\paragraph{Third-party processing and privacy.} Conditioning sources (d) and (e) send individual respondent records to commercial model APIs. Before any run we will confirm that each data-use agreement permits that transfer, send only the variables listed in the grid (no free text, no geography below country), use provider settings that disable retention or training where offered, and never publish respondent-level model inputs or outputs.

\paragraph{Human subjects.} The survey itself involves no new data collection. The protocol is secondary analysis of de-identified human-subjects microdata; we will obtain an exemption determination from an independent IRB, a commercial one if no academic partner's board takes it, named in the registration, before any run rather than assert one.

\paragraph{Harms and dual use.} Simulated populations can misportray and flatten identity groups \citep{wang2024flatten}, and substituting them for participants can stand in for inclusion without delivering it \citep{agnew2024}. Presenting a model's narrative as a person's testimony invites both harms. Simulations can also be used to test persuasive messages on synthetic groups, to rehearse election messaging, or to publish ``opinions'' of populations, often non-Western ones, where this survey finds errors largest. Labeling outputs as model output is necessary but not sufficient. A public system should show validation status and error per population, withhold outputs for populations where no validation exists, and not offer per-person persuasion testing. Our product does not exist yet and has not adopted these recommendations; until it does, they bind no one, including us, and readers should not assume it will follow them.

\section{Conclusion}

Four years of evidence support a narrower claim than ``silicon samples'' first suggested. Language models encode a usable prior over how groups and treatments differ: effect directions and rank orders often agree with human data, though rankings can flip with response format or release, part of the advantage of ranks over levels is an artifact of comparing correlations with absolute errors, and whether US levels agree is contested. In the studies we tabulate, simulations do not reproduce the spread, non-Western levels, or individual answers of real populations without human data per person; joint structure is contested, with signs agreeing more often than magnitudes. Where a non-LLM baseline is fit, it is ahead in \cBaseAhead{} and level in \cBaseLevel{} of \cBase{} studies, mixed in \cBaseMixed{}, and behind only where the LLM is given human data beyond a prompt. The four-axis taxonomy is consistent with about half of the cross-study conflicts we identified, a check with little power to fail; country, language and release remain outside it. The useful role that remains is a prior that real data must correct, validated per population and revalidated per model release. The most important open question, whether simulated populations change the way real ones do after events they have not seen, needs forward-dated tests; the closest published evidence \citep{ashokkumar2026,kozlowski2024,pew2026} is encouraging for effects and directions and discouraging for levels.

\bibliographystyle{plainnat}
\bibliography{refs}

\appendix
\section{Coded Matrix of Included Empirical Studies}
\label{app:matrix}

Codes follow Section~\ref{sec:taxonomy}. A4: none = no reference point; base = non-LLM baseline; ceil = test--retest ceiling; unif = uniform or random; floor = human sampling floor or sample-size equivalence; hum = other human benchmark. Read: full = full text; part = abstract plus a full-text search for the coded comparisons; abs = abstract only. PR: peer-reviewed venue per Crossref, publisher, ACL Anthology, OpenReview or arXiv record (rep = institutional report; ws = workshop; ext = extended abstract). Single coder; ``?'' = not determinable from what we read. Source: \nolinkurl{data/corpus.csv}.

{\scriptsize
\begin{longtable}{@{}L{3.3cm}L{2.4cm}llllll@{}}
\toprule
Study & Country & A1 & A2 & A3 & A4 & Read & PR \\
\midrule
\endhead
% Generated by data/counts.py from corpus.csv. Do not edit.
\citet{argyle2023} & US & d & i? & U2,U3 & none & abs & yes \\
\citet{bisbee2024} & US & c & i & U1,U3 & none & full & yes \\
\citet{santurkar2023} & US & a,b & ii & U1,U2 & hum & abs & yes \\
\citet{durmus2023} & global & a,b & ii & U1 & none & abs & yes \\
\citet{ryan2024} & global & a & ? & U1 & none & abs & yes \\
\citet{jang2026} & US & c & i,iv & U1 & none & part & yes \\
\citet{doudkin2026} & 3 countries & a,c & i,iii,iv & U1,U2 & base & part & --- \\
\citet{park2024} & US & c,e,f & ? & U4 & ceil & abs & --- \\
\citet{ku2026} & Taiwan & e & i & U4 & base & part & yes \\
\citet{peng2025funhouse} & US & a,c,e & i & U1,U4 & floor,unif & full & yes \\
\citet{toubia2025twin} & US & e & i & U4 & ceil,unif & full & yes \\
\citet{ozkan2026} & Turkey & d & i,iii & U1 & none & full & --- \\
\citet{lee2026openweight} & US & a,e & i & U1,U3 & floor & full & --- \\
\citet{tao2024} & 107 countries & a,b & ? & U1,U2 & none & full & yes \\
\citet{wang2024flatten} & US & c & ? & U2 & none & abs & yes \\
\citet{lukauskas2026} & Lithuania & c & i & U1,U3 & base,ceil & part & --- \\
\citet{tornberg2026} & US & a & i & U1 & none & abs & --- \\
\citet{hu2025simbench} & 130+ countries & a,b & ii,iii & U1 & unif & full & yes \\
\citet{suh2025subpop} & US & b & v & U1 & floor,unif & full & yes \\
\citet{zhao2024wvb} & global & c & i & U1 & none & abs & yes \\
\citet{chen2026synthetic} & US, global & c & i & U2,U4 & base & part & --- \\
\citet{meister2024} & US, global & b & i,ii,iii & U1 & unif,hum & full & yes \\
\citet{kambhatla2025} & global & a,b & ii,iii & U1 & none & full & yes \\
\citet{krsteski2025} & ? & c & i,v & U1 & base & part & yes \\
\citet{liu2025global} & global & b & ? & U1 & none & abs & yes \\
\citet{wang2026ess} & 30 countries & b,c & i? & U2,U4 & base & part & --- \\
\citet{williams2026beyond} & global & c & i,v & U1,U3 & none & abs & yes \\
\citet{kim2026korean} & Korea & c & i & U1,U2 & base & abs & yes \\
\citet{horton2023} & --- & c & i & U5 & none & abs & ext \\
\citet{aher2023} & US & c & i & U5 & none & abs & yes \\
\citet{lam2026fast} & Singapore & d? & i & U1 & none & abs & --- \\
\citet{abels2026} & 10 countries & c & ? & U1 & none & full & --- \\
\citet{schroder2025} & --- & a & i & U1 & none & abs & --- \\
\citet{park2023diminished} & US & a & i & U5 & none & abs & yes \\
\citet{dominguez2024} & US & a,c & ii & U1 & unif & abs & yes \\
\citet{binz2025centaur} & --- & e & v & U4 & base & abs & yes \\
\citet{pew2026} & US & d & i & U1 & none & full & rep \\
\citet{bulte2026} & global & b & ? & U1 & none & abs & yes \\
\citet{atari2023} & global & a & ? & U1,U2 & none & abs & --- \\
\citet{ashokkumar2026} & US & a & ? & U5 & hum & abs & yes \\
\citet{rystrom2025} & 4 languages & a & ? & U1 & none & abs & ws \\
\citet{tjuatja2024} & US & a & i & U5 & none & abs & yes \\
\citet{hu2024persona} & --- & c & i & U4 & base & part & yes \\
\citet{beck2024} & --- & c & i & U4 & none & abs & yes \\
\citet{moon2024anthology} & US & c & i & U1 & none & abs & yes \\
\citet{sun2024random} & US & c & i & U1 & none & abs & --- \\
\citet{vonderheyde2025} & Germany & d & i & U1,U2,U4 & none & abs & yes \\
\citet{cao2025} & global & b & v & U1 & none & abs & yes \\
\citet{kim2023augmented} & US & d,e & v & U1,U4 & base & part & --- \\
\citet{sun2023sociodemo} & US & c & i & U4 & none & abs & yes \\
\citet{kozlowski2024} & US & b & ? & U2 & unif & abs & --- \\
\citet{boelaert2025} & ? & ? & ? & U1 & none & abs & yes \\
\bottomrule

\end{longtable}
}

\end{document}
